\chapter{Addiction Rehab}

\section*{Definition and Overview}
Addiction rehab (short for rehabilitation) refers to structured treatment programmes that help people recover from addiction to alcohol or other drugs. Rehab usually involves a period of time away from the person's usual environment, with medical care, counselling, group therapy, and skills training.

Programmes may be residential (where the person stays at the facility) or day-based, and typically last from a few weeks to several months. The goal is not just to stop using, but to build the skills, routines, and support needed to stay well after the programme ends.

\section*{Epidemiology}
Residential rehab is used by a minority of people with substance use disorders, but it is an important part of the treatment system. People are more likely to be offered or to choose rehab when they have moderate-to-severe dependence, have tried outpatient treatment without success, lack a safe home environment, or have other serious health and social problems.

Demand for rehab beds often exceeds supply, and waiting lists are common. Programmes also vary widely in cost, staffing, and philosophy, so careful choice and assessment are important.

\section*{Aetiology and Pathophysiology}
Rehab treats the same underlying brain changes that cause addiction: altered reward circuits, strong cravings, and weakened impulse control. Because these changes persist after the substance is removed, treatment works best when it is intensive and structured, giving the brain and the person time to reset away from everyday triggers.

Many people entering rehab also have withdrawal symptoms, mental health conditions, or physical health problems that must be managed at the same time. A stay in rehab therefore combines a safe, supervised environment with medical care, psychological therapy, and the deliberate building of healthier habits.

\section*{Clinical Features}
Rehab programmes vary, but most share common features:

\begin{itemize}
    \item \textbf{Assessment on admission:} evaluation of substance use, physical health, and mental state
    \item \textbf{Detoxification:} medically supervised withdrawal for those who need it
    \item \textbf{Structured timetable:} a daily schedule of therapy, education, and activities
    \item \textbf{Individual and group therapy:} regular counselling and group sessions
    \item \textbf{Family involvement:} education and sessions for partners and relatives
    \item \textbf{Aftercare planning:} relapse prevention and follow-up arrangements made before leaving
\end{itemize}

\section*{Differential Diagnosis}
Choosing the right level of care matters, and alternatives to rehab include:

\begin{itemize}
    \item Community or home-based withdrawal management (detox)
    \item Outpatient counselling and medication while living at home
    \item Mutual support groups such as Alcoholics Anonymous
    \item Medical or psychiatric treatment without addiction rehabilitation
\end{itemize}

\section*{When to Seek Medical Care}
A person should consider rehab when addiction is having serious effects on health, safety, or daily life, when attempts to stop on their own keep failing, or when professional advice recommends a structured programme. A GP or addiction service can help with assessment and referral.

\textbf{Call 000 (or local emergency services) for:}
\begin{itemize}
    \item Overdose: severe drowsiness, slowed breathing, or unresponsiveness
    \item Seizures, confusion, or delirium during withdrawal
    \item Suicidal thoughts or self-harm
    \item Any life-threatening reaction to drugs or alcohol
\end{itemize}

\section*{Diagnosis and Investigations}
\begin{itemize}
    \item \textbf{Addiction assessment:} a detailed history of substance use and its consequences
    \item \textbf{Medical examination:} to identify withdrawal risk and physical health problems
    \item \textbf{Mental health assessment:} to detect depression, anxiety, or other conditions
    \item \textbf{Blood tests:} liver function, full blood count, and infection screening (hepatitis, HIV)
    \item \textbf{Toxicology:} urine or breath drug testing
    \item \textbf{Risk assessment:} of overdose, self-harm, and pregnancy where relevant
\end{itemize}

\section*{Management}
\begin{itemize}
    \item \textbf{Detoxification:} medically supervised withdrawal, which is essential for alcohol and benzodiazepines
    \item \textbf{Medicines:} to reduce craving or block the effects of drugs, such as buprenorphine or naltrexone
    \item \textbf{Individual therapy:} cognitive behaviour therapy and motivational interviewing
    \item \textbf{Group therapy:} structured groups and 12-step-style meetings
    \item \textbf{Family therapy:} to repair relationships and build support at home
    \item \textbf{Life skills training:} employment, budgeting, communication, and relapse prevention
    \item \textbf{Aftercare:} a clear plan for the weeks and months after leaving rehab
\end{itemize}

\section*{Complications}
\begin{itemize}
    \item Relapse after leaving, particularly in the first year
    \item Withdrawal complications if detoxification is poorly managed
    \item Overdose if someone relapses after a period of reduced tolerance
    \item Leaving the programme before completion
    \item Difficulty readjusting to everyday life and old surroundings
    \item Financial cost and time away from family and work
\end{itemize}

\section*{Prognosis and Outcomes}
Completing a rehab programme is associated with better outcomes than leaving early, but rehab alone does not guarantee lasting recovery. People who finish a programme and continue with aftercare, mutual support groups, and follow-up are more likely to stay well.

Many people who enter rehab relapse at some point and then return to treatment; with each attempt, the chance of lasting recovery improves. Rehab is best seen as one important step within a longer journey of recovery rather than a cure in itself.

\section*{Prevention and Self-Care}
\begin{itemize}
    \item Research programmes in advance and choose one that fits the person's needs
    \item Attend all sessions and engage honestly with therapy
    \item Involve family and close friends in the process
    \item Build a strong aftercare plan before leaving
    \item Keep in touch with recovery groups and support services after discharge
    \item Plan how to handle triggers and high-risk situations
    \item Seek help immediately after a lapse rather than waiting
\end{itemize}

\section*{Sources and Further Reading}
The following organisations provide reliable information on addiction rehabilitation and treatment services.

\begin{itemize}
    \item NHS. \textit{Information on drug and alcohol rehabilitation services.} https://www.nhs.uk/conditions/
    \item World Health Organization. \textit{Resources on substance use treatment and rehabilitation.} https://www.who.int/
    \item NICE. \textit{Clinical guidance on drug misuse and alcohol dependence treatment.} https://www.nice.org.uk/
    \item Mayo Clinic. \textit{Guide to addiction treatment and rehabilitation programmes.} https://www.mayoclinic.org/
    \item MedlinePlus. \textit{Health information on drug and alcohol rehabilitation.} https://medlineplus.gov/
\end{itemize}
